\section{Locality replaces the full tetrahedron count}\label{sec:regions}

The uniform commitment bound still has an exponent linear in $t$, even when
only a small part of a triangulation participates in a useful sequence. The
next result charges the exponent to the \emph{initial} region that the trace
is allowed to consume. Newly created tetrahedra are included automatically
in the incarnation accounting.

\begin{definition}[Allowed initial region]
Let $G_T$ be the simple graph on the original tetrahedra, with adjacency when
they share a paired face. An allowed initial region is a set
$F\subseteq V(G_T)$. A trace is supported in $F$ if no original tetrahedron
outside $F$ is consumed. It may consume every subsequently created tetrahedron.
For $r,c,U\ge0$, let $\Reach_{r,c,U}(T,h)$ be the union of bounded-upward
families supported in regions $F$ with $|F|\le r$ and at most $c$ connected
components in $G_T[F]$.
\end{definition}

We use effective integer caps $0\le r\le t$ and $0\le c\le r$;
larger supplied caps can be reduced to these values without changing the
family. The empty region is included, so zero caps still allow the source
endpoint to be checked.

The region is a \emph{cover} of the actual initial consumption footprint.
Some members of $F$ may survive. An actual consumed set can have several
components while lying inside one connected allowed region. This distinction
matches the executable search and is useful when a small connecting corridor
is included in the supplied region.

\subsection{A supplied region}

\begin{theorem}[Region-sensitive commitment bound]\label{thm:fixed-region}
For a supplied region $F$ of size $i$, the family supported in $F$ with at
most $U$ upward events has at most $11^{3i+5U}$ commuting trace classes.
When $U=0$, the bound is $7^{3i}$. Both the commitment search and the sleep
search exhaust this family with polynomial overhead per encoding or
representative, apart from endpoint-predicate work.
\end{theorem}
\begin{proof}
Give every original tetrahedron outside $F$ the forced symbol idle. Such
tetrahedra are protected forever. After $u$ upward and $d$ downward events,
at least $t-i$ original tetrahedra remain, so
\[
 t+u-d\ge t-i,\qquad d\le i+u.
\]
Only the $i$ initially eligible tetrahedra and all newly born tetrahedra need
freely chosen commitments. Their number is
\begin{equation}\label{eq:active-incarnations}
 i+3u+2d\le3i+5u\le3i+5U.
\end{equation}
The proof of Theorem~\ref{thm:commitment} preserves the identities of consumed
original tetrahedra under every interchange, so it remains valid with the
protection restriction. The sleep-set proof also preserves this restriction.
Every generated trace is supported in $F$, and every trace supported in $F$
has its prescribed idle commitments on the protected originals.
\end{proof}

The number of events in such a trace is at most $i+2U$, even when the
surrounding triangulation is much larger. The code checks the two inequalities
in this proof at every search frame. A bad bookkeeping update therefore
cannot silently invalidate the claimed measure.

\subsection{Counting and enumerating regions}

The maximum number of distinct neighbours of a vertex of $G_T$ is four.
Loops and repeated face adjacencies do not increase this bound.

\begin{lemma}[Counting connected regions]\label{lem:connected-count}
In a graph with $t\ge1$ vertices and maximum degree at most four, the number
of connected $i$-vertex subsets is at most $t16^{i-1}$ for $i\ge1$.
The number of subsets of size at most $r$ with at most $c$ connected
components, including the empty set, is at most
\begin{equation}\label{eq:region-count}
 (r+1)(c+1)t^c32^r.
\end{equation}
\end{lemma}
\begin{proof}
Root a connected set at its least-labelled vertex. Choose a deterministic
spanning tree of the induced graph and traverse its edges depth first,
each twice. Record the root and the chosen local neighbour port at each
of the $2(i-1)$ steps. There are at most four choices at each step. The
walk reconstructs the visited vertex set, so the encoding is injective
on the connected sets selected by the deterministic rule. This proves
$t4^{2(i-1)}$.

For a set of size $i$ with $j\ge1$ components, order the components by their
least vertex. Encode their $j$ roots, their positive sizes summing to $i$,
and their walks. This gives the upper bound
\[
 t^j\binom{i-1}{j-1}16^{i-j}\le t^j32^i.
\]
Summing over $1\le i\le r$ and $1\le j\le\min(c,i)$, then adding the empty
set, proves the intentionally coarse uniform bound
\eqref{eq:region-count}. It also holds for $r=0$ or $c=0$.
\end{proof}

Counting alone would not give a useful algorithm if finding those regions
required enumerating all $t^r$ subsets. The artifact uses the following
explicit reverse search.

\begin{lemma}[Canonical parent and complete reverse search]\label{lem:reverse}
For a connected set $A$ with root $\rho=\min A$ and $|A|>1$, delete the
largest $x\in A\setminus\{\rho\}$ for which $G_T[A\setminus\{x\}]$ stays
connected. This defines a unique parent. Starting from every singleton,
and accepting an added frontier vertex precisely when this rule returns
the current set as parent, enumerates each connected region once, with
polynomial overhead in its size.
\end{lemma}
\begin{proof}
A spanning tree of $G_T[A]$ has a leaf other than $\rho$. Deleting that leaf
leaves a spanning tree on the other vertices, hence preserves connectivity
of the induced graph. At least one eligible deletion therefore exists.
The largest-label rule makes it unique. Iterating the parent decreases
size and preserves the root, eventually reaching $\{\rho\}$.

Any child of $A$ adds a vertex adjacent to $A$: otherwise it is not connected.
The algorithm therefore considers it among the frontier vertices and accepts
it at its unique parent. It cannot be accepted at another parent. There are
at most $4|A|$ distinct frontier candidates; each parent computation needs
at most $|A|$ connectivity tests on at most $|A|+1$ vertices. This is
polynomial overhead, with no enumeration of vertex-addition permutations.
\end{proof}

For several components, the code selects connected pieces in increasing
order of their least vertex. After choosing a component, it excludes its
closed neighbourhood from later components. The pieces are consequently
the actual components of the permitted region. Every allowed region has
exactly one such ordered decomposition. The same reverse-search argument
works in the remaining induced graph. An explicit stack avoids a recursion
limit tied to the number of vertices.

\subsection{The constructive locality theorem}

\begin{theorem}[Complete local reachability]\label{thm:local}
Let $(T,h)$ be a valid source with $t\ge1$ tetrahedra and initial cocycle
bit bound $B$. Suppose a chosen endpoint predicate is invariant under the
tracked marked isomorphisms and has polynomial bit complexity. Then
$\Reach_{r,c,U}(T,h)$ can be searched exhaustively in
\begin{equation}\label{eq:local-bound}
 t^c2^{O(r+U)}\poly(t+U+B)
\end{equation}
bit operations. A precise sufficient counting factor is
\[
 (r+1)(c+1)t^c32^r11^{3r+5U}.
\]
Every successful endpoint is geometrically replayable from the supplied
source. If the endpoint predicate has a different worst-case cost, that
cost must be included as an additional factor.
\end{theorem}
\begin{proof}
Enumerate the regions by Lemma~\ref{lem:reverse} and its component extension.
For each region use Theorem~\ref{thm:fixed-region}. The union is exactly the
defined family: membership requires containment in at least one such region.
Repeated endpoints across different regions do not affect correctness or
the upper bound. Lemma~\ref{lem:connected-count} bounds the number of regions.
The polynomial marking and geometric costs are established in
Section~\ref{sec:verification}. Multiplying these factors proves the claim.
\end{proof}

\begin{corollary}[Polynomial and quasi-polynomial classes]\label{cor:quasi}
For one connected allowed region, $r+U=O(\log t)$ gives a polynomial search
bound, and $r+U=O(\log^2 t)$ gives $t^{O(\log t)}\poly(B)$.
More generally, $c=O(\log t)$ together with $r+U=O(\log^2 t)$ gives the
latter quasi-polynomial bound.
\end{corollary}
\begin{proof}
Take logarithms of the non-polynomial factor in
\eqref{eq:local-bound}: it is $O(c\log t+r+U)$. Substitute the stated
conditions. The fixed polynomial factors do not change these classifications.
\end{proof}

This improves what a direct move-string analysis provides for the same
locality conditions. A trace has at most $r+2U$ events but can choose among
$O(t+U)$ locations at each event. That direct bound pays
$O((r+U)\log(t+U))$ in the exponent. The region theorem pays
$O(c\log t+r+U)$, isolating the cost of choosing the spatial components.
The improvement concerns the explicit search model, not a lower bound on
every alternative way to recognize knots.

\subsection{A precise remaining global obligation}

Suppose an outer, source-certified algorithm makes $Q(n)$ calls of this
kind on sources with parameters $t_i,B_i,r_i,c_i,U_i$. Its accumulated cost
is bounded by the sum of the individual costs, plus its other work:
\begin{equation}\label{eq:outer-sum}
 \sum_{i=1}^{Q(n)}t_i^{c_i}2^{O(r_i+U_i)}\poly(t_i+U_i+B_i).
\end{equation}
Thus polynomially bounded source sizes and markings, a quasi-polynomial
number of outer calls, and the bounds in Corollary~\ref{cor:quasi} would
yield a quasi-polynomial total. This is a sufficient accounting statement.
It does not prove completeness of such an outer algorithm or bound its
restarts. In particular, the existence of a short successful trace for an
unknot alone says nothing about how a negative instance is decided.

The outstanding geometric target is now concrete: establish a useful
small-region coverage theorem for a specified source construction and
surface family, including its behavior after certified reductions. The
existing coherent-height obstruction shows that coverage cannot be asserted
without allowing changes of triangulation or a richer surface representation.
